\chapter{Introduction}
\label{ch:intro}

\section{From Pilot to Production}
\label{sec:intro-pilot}

The v1.0 report of this project established a hypothesis: that 3D Gaussian Splatting (\acrshort{3dgs}) could serve as a practical medium for preserving the spatial and experiential qualities of cultural heritage spaces that conventional photography and video cannot capture. That report demonstrated viability. It produced a photorealistic splat from a video walkthrough, observed that mesh extraction degraded fidelity, and recommended that splats be treated as preservation artefacts in their own right.

This report supersedes that pilot. In the intervening period the system was re-platformed onto LichtFeld~Studio~v0.5.2 \parencite{LichtFeldStudio2026}, and the segmentation, generative-reconstruction, and scene-assembly subsystems that the v1.0 report described as future work were completed and integrated. The \texttt{gaussian-toolkit} now realises the project's original design goals in full. From a single hand-held 4K video of a gallery interior, the pipeline produces not merely a splat, but a \emph{structured, editable, hierarchical reconstruction}: a photoreal Gaussian field, a watertight polygonal mesh of the environment, an automatic decomposition of the scene into individually-addressable objects, an independently reconstructed and fully-textured 360\textdegree\ mesh for each of those objects, and a single Universal Scene Description (\acrshort{usd}) scene graph that places every reconstructed object at its surveyed position within the meshed room.

\begin{keyfinding}
The v2.0 system resolves the central tension of the v1.0 pilot---that splat fidelity and mesh editability were mutually exclusive. It maintains the Gaussian splat and the polygonal mesh as co-registered variants of the same scene-graph prims. One capture yields both a photoreal browser experience and a game-engine-ready mesh hierarchy, bound together by surveyed geometry.
\end{keyfinding}

\section{The Preservation Problem}
\label{sec:intro-problem}

Immersive and experiential cultural works---installation art, site-specific performance, interactive exhibitions---resist conventional documentation. Their meaning is inseparable from the act of inhabiting them: the spatial relationships between objects, the way light falls across a room, the sequence in which a visitor encounters a work. A photograph flattens this; a video fixes a single path through it. Neither preserves the \emph{navigable space} itself \parencite{UNESCO2003,Champion2011}.

Volumetric reconstruction offers a different proposition: capture the space as a navigable three-dimensional field, from which any viewpoint can be synthesised after the fact. The challenge has historically been one of cost and expertise---laser scanning and photogrammetry demand specialist equipment, controlled conditions, and skilled operators \parencite{Remondino2011,Stylianidis2016}. The contribution of this project is to collapse that barrier: a single uncalibrated video, captured on a consumer phone, processed by an automated pipeline, yielding a result that is simultaneously archival-grade and audience-ready.

\section{Design Goals and Their Achievement}
\label{sec:intro-goals}

The project was funded by the University of Salford to trial whether contemporary AI reconstruction could meet a specific set of capability goals. Table~\ref{tab:intro-goals} states each goal as originally framed and the v2.0 outcome against it.

\begin{table}[htbp]
  \centering
  \caption{Stated design goals and their realisation in v2.0.}
  \label{tab:intro-goals}
  \begin{tabularx}{\textwidth}{@{}L{5.4cm}X@{}}
    \toprule
    \textbf{Design goal} & \textbf{v2.0 outcome} \\
    \midrule
    Photoreal capture of a heritage interior from consumer video &
    Achieved. A 4K phone walkthrough yields a 1.18M-Gaussian field at \acrshort{psnr}~29.6~dB in $\sim$24~minutes of \acrshort{gpu} time (Chapter~\ref{ch:reconstruction}). \\
    \addlinespace
    A polygonal mesh of the environment with a hierarchical scene graph &
    Achieved. Four pluggable backends extract the room mesh (Chapter~\ref{ch:meshing}); the \acrshort{usd} scene graph organises the result hierarchically (Chapter~\ref{ch:scenegraph}). \\
    \addlinespace
    Separate, 360\textdegree-reconstructed objects within the scene &
    Achieved. Salient objects are reconstructed independently as watertight, textured meshes via multi-view generative inference (Chapter~\ref{ch:scenegraph}). \\
    \addlinespace
    Objects extracted automatically by segmentation &
    Achieved. \acrshort{sam}2 decomposes the scene into 31 coherent object clusters with 98.3\% Gaussian coverage (Chapter~\ref{ch:scenegraph}). \\
    \addlinespace
    Objects turned into textured polygonal meshes by a generative engine &
    Achieved. Per-object orbit renders drive Hunyuan3D~2.0 through ComfyUI, with FLUX.1 background completion, yielding \acrshort{pbr}-textured meshes (Chapter~\ref{ch:scenegraph}). \\
    \addlinespace
    Emplacement of objects into the polygonal reconstruction &
    Achieved. Each object mesh is positioned at its \acrshort{colmap}-surveyed centroid within the environment mesh, in one scene graph (Chapter~\ref{ch:scenegraph}). \\
    \bottomrule
  \end{tabularx}
\end{table}

\section{Contributions of v2.0}
\label{sec:intro-contributions}

The technical contributions documented in this report are:

\begin{enumerate}
  \item \textbf{A dual splat-and-mesh representation under a single scene graph.} The system co-registers the photoreal Gaussian field and the polygonal mesh as \acrshort{usd} variant sets, so the same artefact serves browser, engine, and archive without re-processing (Chapters~\ref{ch:scenegraph}, \ref{ch:delivery}).

  \item \textbf{An eight-stage agentic pipeline} spanning frame extraction, Fibonacci-sphere keyframe selection, Structure-from-Motion, splat training, segmentation, generative per-object reconstruction, environment meshing, and scene assembly, orchestrated across a three-container compute fabric (Chapter~\ref{ch:pipeline}).

  \item \textbf{Integration of LichtFeld~Studio~v0.5.2}, exploiting its scene-graph harmonisation, native \acrshort{usd} I/O, in-application evaluation, and \acrshort{mcp} control surface to drive training and export programmatically (Chapter~\ref{ch:lichtfeld}).

  \item \textbf{A per-object generative reconstruction subsystem} that completes the unseen geometry of each segmented object---the back of a sculpture a single walkthrough never observes---using multi-view diffusion (Chapter~\ref{ch:scenegraph}).

  \item \textbf{A measured four-way comparison of mesh-extraction backends}---TSDF, MILo, CoMe, and Gaussian wrapping---characterising their runtime, face count, geometric accuracy, and licence posture on the same capture (Chapter~\ref{ch:meshing}).
\end{enumerate}

\section{Report Structure}
\label{sec:intro-structure}

The report is organised in three parts. \textbf{Part~\ref{part:architecture}} establishes the system: the technology foundations (Chapter~\ref{ch:foundations}), the pipeline and container architecture (Chapter~\ref{ch:pipeline}), and the LichtFeld~Studio integration (Chapter~\ref{ch:lichtfeld}). \textbf{Part~\ref{part:results}} presents the reconstruction results: the end-to-end 4K case study (Chapter~\ref{ch:reconstruction}), the hierarchical scene graph and per-object emplacement that constitute the project's headline achievement (Chapter~\ref{ch:scenegraph}), and the mesh-backend comparison (Chapter~\ref{ch:meshing}). \textbf{Part~\ref{part:delivery}} covers delivery to browser, engine, and archive (Chapter~\ref{ch:delivery}) and the deployment roadmap for the University of Salford (Chapter~\ref{ch:roadmap}).

\begin{technote}
All quantitative results in this report derive from a single real 4K capture of a gallery interior, processed end-to-end on the project's dual-GPU workstation. Numbers are consistent across chapters: where a metric appears more than once, it refers to the same processing run.
\end{technote}
